% Chapter 1 -- The Apple Menu (What's New in System 7 "Apple Menu Items"
% topic form, apple-macintosh.md Sec 8.3).  Sources (Law 1):
% os/flair/apple_menu.c and spec/apple_menu_items.h (About This
% Computer..., a divider, then each accessory whose program is installed,
% in alphabetical order: "omit, never gray"); os/milton/kmain.c
% (flair_live_apple_installed asks the disk each time the menu is pulled
% down; flair_live_about: BANNER_LINE1, Built-in Memory = 1024K + the
% extended memory stage2 found, Largest Unused Block from the InitechDOS
% memory chain; APPLE-LAUNCH through the same launch verb as a
% double-click).  The menu itself is registered in the FLAIR User's Guide,
% Chapter 8.
% GATES: test-apple-menu test-flair-apple-menu test-flair-calcda test-pubs-flair-screens
\Gates{test-apple-menu, test-flair-apple-menu, test-flair-calcda,
  test-pubs-flair-screens}
\chapter{The Apple Menu}

The Apple menu, at the left end of the upper menu bar, is available
whatever you are doing.  It tells you about your computer, and it lists the
accessories on the data disk, so that you can open one without going back
to the Finder.

\FlairScreen[scale=4, label=fig:applemenu]{applemenu}{The Apple menu}

\section{About This Computer}
\index{About This Computer}
\Operation{APPLE-ABOUT}
Choose \menuitem{About This Computer...} to see an alert box with three
lines of information.

\FlairScreen[label=fig:about]{about}{About This Computer}

\SideHead{Version}The first line names the operating system on the
startup disk and its version: \texttt{InitechDOS Version 3.30}.  Quote it
when you call the Help Desk.

\SideHead{Built-in Memory}All the memory of the computer, in kilobytes:
the first 1024K and the extended memory above it.

\SideHead{Largest Unused Block}The largest piece of the memory InitechDOS
gives to its programs that is free at this moment.  When FLAIR has not
handed that memory to InitechDOS, the line reads \texttt{none}.

Click \texttt{OK} or press \key{Enter} to put the alert away.

\section{Accessories}
\index{accessories}
Below the line, the Apple menu lists one command for each accessory whose
program is on the data disk, in alphabetical order.  The data disk as
delivered holds one accessory, the Calculator (Chapter~2), whose program is
\file{CALCULAT.EXE} in the \file{APPS} folder.

FLAIR looks at the disk each time you pull the menu down.  If an
accessory's program is not in its folder, the accessory is simply not
listed --- FLAIR never shows a command you cannot use.  If the menu holds
no accessories, the line is left out too.

Choosing an accessory opens it just as double-clicking its icon would.  An
accessory is an application from the disk, so it does not open while
another one is open: quit the other application first.
